\chapter{Metodología y arquitectura del banco de pruebas}
\label{cap:metodologia}

\section{Arquitectura general}

El banco de pruebas se organiza en seis fases encadenadas, cada una cerrada por un hito
(\emph{gate}) con criterio de aceptación medible:

\begin{enumerate}
  \item[\textbf{F1.}] \textbf{Gemelo termodinámico} del CFM56-7B26 en pyCycle: ciclo de diseño,
    fuera de diseño, \emph{decks} de línea base sobre la envolvente y matriz de coeficientes de
    influencia (ICM). \emph{Gate} H1: error $\leq 3$\,\% frente a los datos certificados en los
    puntos ancla.
  \item[\textbf{F2.}] \textbf{Generador de flota sintética} (SynCFM56): trayectorias de salud por
    mecanismo, flota jerárquica de $\sim$100 motores \emph{run-to-failure}, modelo de sensores
    (ruido, sesgo, cuantización, huecos) y \emph{snapshots} tipo ACARS con verdad-terreno.
    \emph{Gate} H2: dataset auditado, sin fuga entre particiones, con dificultad verificada.
  \item[\textbf{F3.}] \textbf{EHM tradicional} de referencia: línea base y desviaciones
    suavizadas, \emph{trending} con reglas de persistencia, GPA-WLS
    (\cref{eq:wls}), Kalman-GPA (\cref{eq:kalman}), aislamiento por reglas expertas y RUL por
    extrapolación robusta del margen de EGT. \emph{Gate} H3: métricas completas sobre test.
  \item[\textbf{F4.}] \textbf{EHM con IA}: detección de anomalías (autoencoder, Isolation
    Forest), diagnóstico multiclase (XGBoost, redes), pronóstico de RUL (LSTM/TCN/Transformer),
    híbrido \emph{physics-informed}, incertidumbre \emph{conformal} y explicabilidad
    física (PCS). \emph{Gate} H4.
  \item[\textbf{F5.}] \textbf{Evaluación comparativa pre-registrada}: métricas comunes,
    contrastes estadísticos, ablaciones y validación \emph{sim-to-real} en N-CMAPSS.
    \emph{Gate} H5.
  \item[\textbf{F6.}] Casos de estudio, \emph{dashboard} y redacción. \emph{Gate} H6.
\end{enumerate}

\section{Salvaguardas metodológicas}
\label{sec:protocolo}

La credibilidad de una comparación IA-contra-tradicional depende de eliminar los sesgos que
plagan la literatura. Se adoptan cuatro salvaguardas:

\begin{description}
  \item[Pre-registro.] Hipótesis (\cref{tab:hipotesis}), métricas, particiones (con \emph{hash}
    de los ficheros), contrastes y regla de parada se congelan en el repositorio con una
    etiqueta de git (\codigo{prereg-v1}) \emph{antes} de evaluar sobre el conjunto de test.
  \item[Presupuesto de ajuste simétrico.] Los hiperparámetros del pipeline tradicional
    (umbrales, $\lambda$, $\mathbf{Q}$, ventanas) y de cada familia de IA se ajustan con el
    mismo número de \emph{trials} de Optuna, usando solo entrenamiento y validación.
  \item[Partición por motor.] Ningún motor aparece en más de una partición
    (70/10/20 motores en train/val/test), con auditoría automática anti-fuga en integración
    continua.
  \item[Baseline fuerte.] El pipeline tradicional se implementa con el mismo cuidado que la IA,
    documentando cada regla con su justificación física; el \emph{smearing} se mide contra la
    verdad-terreno, no se asume.
\end{description}

\section{Reproducibilidad e infraestructura}
\label{sec:reproducibilidad}

Todo el proyecto es un repositorio público\footnote{\url{https://github.com/cedarcreeks/EHMbrAIn}}
con entorno reproducible: Python~3.11 gestionado con \codigo{uv} y versiones bloqueadas
(\codigo{uv.lock}); pyCycle~4.4 sobre OpenMDAO~\cite{Gray2019}; configuraciones declarativas en
YAML; y pruebas automatizadas (pytest) ejecutadas en integración continua en cada cambio. Los
datos y experimentos de fases posteriores se versionarán con DVC y MLflow respectivamente.

Dos decisiones concretas ilustran el estándar de rigor perseguido:

\begin{itemize}
  \item \textbf{Contrato de calibración versionado.} Los objetivos numéricos del modelo
    (\cref{tab:tcds}) viven en un único fichero
    (\codigo{conf/cfm56\_7b\_targets.yaml}) con valor, tolerancia, fuente y estado
    (\codigo{VERIFIED}/\codigo{TO\_VERIFY}) por entrada. Ese fichero lo leen tanto los
    ejecutores del modelo como las pruebas: un cambio de objetivo es un cambio visible en el
    control de versiones, nunca una edición silenciosa.
  \item \textbf{Pruebas canario de dependencias.} Durante la puesta en marcha se detectó una
    incompatibilidad entre pyCycle~4.4 y numpy~$\geq 2$ (error en el cálculo termodinámico
    tabular). Se fijó \codigo{numpy<2} y se dejó una prueba de humo que ejecuta un ciclo
    completo en cada ejecución de CI, de modo que cualquier actualización de dependencias que
    rompa el solver falle de forma visible e inmediata.
  \item \textbf{Cifras autogeneradas.} Todas las figuras y tablas de resultados de este
    documento las genera un script (\codigo{scripts/make\_memoria\_assets.py}) que resuelve el
    modelo y escribe directamente los PDF y los ficheros \LaTeX{}
    (\codigo{paper/memoria/tablas/*.tex}), junto con un \codigo{assets.json} de
    trazabilidad. Ningún número de resultados se copia a mano.
\end{itemize}

\section{Métricas de comparación}

Ambas familias se evalúan con métricas idénticas: para detección, \emph{lead time}, tasa de
falsas alarmas por mil vuelos, F1 y PR-AUC; para aislamiento, exactitud top-1/top-2, macro-F1,
índice de \emph{smearing} (fracción de $\|\hat{\mathbf{x}}\|_1$ asignada a componentes sanos) y
PCS; para la estimación de salud, RMSE de $\hat{\mathbf{x}}$ frente a la verdad-terreno; para el
RUL, RMSE, MAE, \emph{score} NASA PHM08
\begin{equation}
  S=\sum_i s_i,\qquad
  s_i = \begin{cases} e^{-d_i/13}-1 & d_i<0\\ e^{d_i/10}-1 & d_i\geq 0\end{cases},
  \qquad d_i = \widehat{\mathrm{RUL}}_i - \mathrm{RUL}_i,
\end{equation}
exactitud $\alpha$--$\lambda$ y horizonte de pronóstico; y para la incertidumbre, cobertura
empírica frente a nominal y anchura media de los intervalos.
